\section{SM occupancy and register pressure}
\label{sec:sm_occupancy}

\subsection{The \texttt{nuc\_t[MAX\_NUC\_PER\_MAT]} stack array}

The hottest CUDA kernel, \texttt{trace\_and\_sample}, walks each
particle to its next interaction. The macroscopic cross section
at the particle's energy is the sum of every microscopic cross
section in the material, weighted by atom density:
\[
\Sigma_t(E)
= \sum_{i \in \text{material}} N_i \,\sigma_t^{(i)}(E).
\]
The kernel evaluates every term as it computes the sum, and
\emph{also} retains each $N_i \sigma_t^{(i)}(E)$ in a per-thread
stack array \texttt{nuc\_t[MAX\_NUC\_PER\_MAT]} because the
reaction-selection sampler that runs after the sum needs to draw a
nuclide proportional to the per-nuclide contribution. The
straightforward implementation stores all of those contributions,
then samples on a second pass.

\texttt{MAX\_NUC\_PER\_MAT} is the single source-of-truth constant
\texttt{crate::MAX\_NUCLIDES\_PER\_MATERIAL = 128} that we bumped
from 32 last year to accommodate HMF-069 (69 nuclides) and Pu
solution cases (67). The CPU side imports the constant; the GPU
side receives it via NVRTC \texttt{-DMAX\_NUC\_PER\_MAT=N} so
\texttt{transport.cu} sizes the stack array correctly. The header
\texttt{\#error}s out if NVRTC forgets the flag --- silent
breakage was an early hazard.

The cost: $128 \times 8\,\text{B} = 1\,\text{KB}$ of per-thread
stack on Ampere. On sm\_86 each SM has 65\,536 32-bit registers
total; with $1\,\text{KB}$ of stack array allocated, plus the
other live registers in the trace-and-sample body (energy bin
index, sampled cumulative, RNG state, particle vector, surface /
cell descent stacks), the compiler often promotes the array
into stack-frame memory (local memory in CUDA terminology), which
serializes to L1 / L2 instead of staying in registers.

We considered streaming the per-nuclide loop, computing each
$N_i \sigma_t^{(i)}(E)$ and immediately folding it into a running
cumulative without retaining \texttt{nuc\_t[]}, then evaluating
the cross sections a second time during reaction selection. This
doubles the XS evaluation cost but eliminates the register
spill. We did not commit to either branch: the trade-off depends
heavily on whether the per-nuclide loop is bound on register
pressure (favouring streaming) or on XS evaluation cost
(favouring caching). Both are documented in the source comments
as future work and are part of the agenda when we get a
representative sweep on H100-class hardware where the larger
register file changes the trade-off.

\subsection{Compute-capability-specific tuning constants}

A handful of magic numbers in the kernel-launch heuristics depend
on the device:

\begin{itemize}
    \item \textbf{Blocks per SM.} sm\_86 (Ampere RTX A1000 / 3080)
        and sm\_89 (Ada L40 / RTX 4090) limit to 16 blocks per SM;
        sm\_90+ (Hopper, Blackwell) raises this to 32. Our
        \texttt{recommend\_refill\_factor} function reads the
        compute capability via cudarc's
        \texttt{CU\_DEVICE\_ATTRIBUTE\_COMPUTE\_CAPABILITY\_MAJOR}
        and routes both branches.
    \item \textbf{Register file size.} sm\_75 (Turing T4), sm\_86,
        sm\_90 all carry 65\,536 registers per SM. We use a
        conservative ceiling that does not need to fork on this
        constant.
    \item \textbf{atomicAdd(double*, double) availability.}
        Introduced in CC 6.0 (Pascal GP100). We reject earlier
        compute capabilities explicitly at NVRTC compile time
        with a named error message --- silent fallback would have
        meant the kernel compiles but produces wrong tallies on
        any device that synthesises the atomic via lock-cmpxchg.
\end{itemize}

\subsection{What we believed vs.\ what we measured}

Two priors we revised after measurement:

\paragraph{The SM-coverage threshold for refill is geometry-dependent.}
Our first refill-pool experiment assumed that any time the
batch's particle population dropped below the
``$\text{SM count} \times \text{max threads per SM}$'' threshold,
refilling would help. Measurement (§\ref{sec:refill}) showed the
threshold is not a single number: under-occupied initial
populations don't benefit (the problem is at step 0, not in the
tail), and over-occupied initial populations also don't benefit
(no SM headroom for refilled work). Refill wins in the mid-curve
where the initial population fills the SMs but the population
decays over $\sim$5--10 events. The right tuning constant is the
saturation \emph{knee} of the device, not the SM count.

\paragraph{Memory-bandwidth limits do not show up where we expected.}
We expected SVD reconstruction to be compute-bound (one rank-$k$
FMA chain per lookup). Profiling showed the bottleneck is not
the FMA but the divergent loads --- $32$ threads in a warp each
indexing into per-nuclide \texttt{basis} and
\texttt{coeffs} arrays at different offsets and energy bins, a
pattern the Monte-Carlo XS literature catalogues as the
classical memory-bottleneck regime~\cite{tramm2014xsbench}.
The basis and coefficient arrays are cache-blocked (per-nuclide
pointer arrays since \texttt{gpu\_per\_nuclide.rs}'s pointer-table
work), but the \emph{access pattern} across a warp is the
limiter, not the arrays' aggregate bandwidth.
